\chapter*{Conclusiones }


\section*{Conclusiones}

A lo largo del desarrollo de este proyecto se logró consolidar la integración y caracterización de los distintos subsistemas que conforman el generador de \textit{chirp} para la arquitectura de radar FMCW. El estudio y la validación de cada bloque funcional desde la etapa de alimentación y filtrado hasta la síntesis de frecuencia mediante el PLL con la placa de evaluación EV-ADF4159EB3Z y el procesamiento de microondas permitieron contrastar los resultados experimentales obtenidos en el banco de pruebas con los modelos teóricos y las simulaciones previas.

Con el propósito de evaluar el comportamiento dinámico del conjunto y verificar el funcionamiento global del sistema integrado, se llevaron a cabo mediciones específicas orientadas a contrastar la respuesta integral de todos los subsistemas interconectados frente a los requerimientos de operación del radar FMCW.

La implementación del filtro Hairpin demostró ser una solución adecuada para eliminar las componentes de frecuencia indeseadas asociadas al arranque del sistema y atenuar las armónicas generadas por el VCO, logrando un comportamiento acorde a los requerimientos pese a las tolerancias inherentes al proceso de fabricación sobre sustrato FR-4. Asimismo, el divisor de potencia Wilkinson validó una correcta distribución de la señal con niveles de adaptación y aislamiento satisfactorios en la banda de microondas.

Con el fin de dotar al equipo de versatilidad en condiciones de campo y validar su autonomía operativa, el sistema fue diseñado e instrumentado con dos modos de uso diferenciados, permitiendo realizar mediciones específicas incluso en entornos donde no se disponga de fuentes de alimentación conectadas a la red eléctrica tradicional.

Con el objetivo de optimizar la robustez física, facilitar el transporte y garantizar la integridad operativa del conjunto, se diseñó y fabricó una carcasa protectora a medida para el prototipo. Esto permitió mejorar sustancialmente la disposición de los módulos, asegurar una adecuada sujeción mecánica y brindar mayor seguridad frente a manipulaciones externas durante la realización de ensayos y mediciones fuera del laboratorio.

En lo referente al núcleo de generación de señal, la caracterización del oscilador a lazo abierto y su posterior integración dentro del lazo de control cerrado del sintetizador pusieron de manifiesto las ventajas fundamentales de la arquitectura elegida. Si bien el VCO opera de manera autónoma cubriendo el rango de $1550~\mathrm{MHz}$ a $1850~\mathrm{MHz}$, la implementación del PLL resulta indispensable para garantizar la estabilidad a largo plazo, reducir el ruido de fase y asegurar la repetitividad temporal exigida en un sistema de radar FMCW coherente.

La evaluación del sistema bajo diferentes perfiles y tiempos de barrido permitió determinar los límites prácticos de seguimiento del lazo cerrado. Para tiempos de barrido moderados ($0{,}5~\mathrm{ms}$), el sistema sigue fielmente el perfil programado con una desviación mínima. Si bien velocidades de modulación extremas ($0{,}02~\mathrm{ms}$) evidencian las limitaciones dinámicas del lazo, el análisis global confirma que el diseño propuesto cumple satisfactoriamente con los objetivos planteados, proporcionando una plataforma robusta, flexible y escalable para aplicaciones de radiofrecuencia y radar.

\section*{Trabajos a futuro}

A partir de los resultados y las experiencias obtenidas en el transcurso del proyecto, se identifican diversas líneas de mejora y continuidad orientadas a optimizar el desempeño global del sistema:

\begin{itemize}
    \item \textbf{Optimización de los sustratos y técnicas de fabricación:} Migrar del sustrato FR-4 a materiales de alta frecuencia y baja tangente de pérdidas, lo que permitiría mejorar la precisión geométrica de las líneas microstrip, reducir las discontinuidades por transiciones abruptas y optimizar la adaptación de impedancias en los filtros y divisores.
    
    \item \textbf{Integración física modular avanzada:} Desarrollar una placa única que centralice todos los subsistemas con el fin de minimizar las pérdidas de inserción asociadas al uso de cableado coaxial y conectores SMA externos, incrementando la robustez mecánica y la inmunidad al ruido electromagnético del conjunto.
    \item \textbf{Integración y prueba en un radar FMCW:} Desarrollar todo el conjunto de bloques que corresponden a un radar FMCW y detectar objetivos enviando la señal que genera el sistema.
\end{itemize}